\section{Evaluation}
\label{sec:evaluation}

To answer RQ1--RQ3, we study a representative P4 property, ill-typed rejection, which states that ill-typed programs are rejected; for RQ4, we apply \texttt{spec4pbt} to several properties.

\subsection{Setup and Validation}
\label{sec:setup}
\label{sec:validation}

We describe target implementations, agent settings, and budgets. Also, we validate the translation using \texttt{p4c} test cases.

\subsection{RQ1: Feasibility}
\label{sec:rq1}

We assess whether an agent can complete both constructions for ill-typed rejection, reporting completion, time, token cost, and human assistance for each.

\subsubsection{Proving the Property}
\label{sec:rq1-proof}

We ask the agent to prove ill-typed rejection over the P4 specification.

\subsubsection{Constructing and Verifying the Generator}
\label{sec:rq1-generator}

We ask the agent to construct a sound and complete generator of ill-typed programs, comparing runs with and without decomposition support.

\subsection{RQ2: Hallucinations Caught by Lean}
\label{sec:rq2}

We classify the hallucinations that Lean rejects in RQ1, proofs and generators the agent presents as correct, by the false alarms or missed inputs each would have caused in testing.

\subsection{RQ3: Generator Efficiency}
\label{sec:rq3}

We compare the generator's throughput and coverage over time with Failing with Purpose~\cite{lee26fwp}, Palamedes~\cite{goldstein26palamedes}, and QuickChick~\cite{lampropoulos18generators}.

\subsection{RQ4: Bug Finding}
\label{sec:rq4}

We apply \texttt{spec4pbt} to several properties, including ones we write about recent P4 language features.
